\subsection{The subcategory generated by adjacent arrows}

\begin{lemma}[Proved: the adjacent arrows determine top homology]
\label{A7:adjacent-generated-top-homology}
Let $\mathcal C$ be a finite acyclic category with full support
$0<1<2<3<4$.  Let $\mathcal D\subseteq\mathcal C$ be a
subcategory containing every object and every arrow in the four
adjacent hom-sets $\mathcal C(i,i+1)$.  For every field $k$,
the inclusion induces isomorphisms
\[
 H_4(B\mathcal D;k)\xrightarrow{\cong}H_4(B\mathcal C;k),
 \qquad
 H^4(B\mathcal C;k)\xrightarrow{\cong}H^4(B\mathcal D;k).
\]
Consequently, if all degree-two products on $B\mathcal D$ vanish
over $k$, then all degree-two products on $B\mathcal C$ vanish.
In particular, one may take $\mathcal D$ to be the subcategory
generated by the four adjacent hom-sets.
\end{lemma}
\begin{proof}
Every nondegenerate four-simplex has the object string $01234$.
Its four arrows are therefore adjacent arrows, all of which
belong to $\mathcal D$.  Thus the inclusion identifies the
normalized fourth chain groups.  Every face of such a simplex
also belongs to $\mathcal D$, since this is a subcategory.
The inclusion of third chain groups is an injection on their
literal simplex bases and intertwines the fourth boundaries.
It follows that the two fourth-boundary kernels coincide.
Both categories have no nondegenerate simplices of degree five,
so these kernels are their fourth homology groups.  Duality
over the field $k$ gives the stated fourth-cohomology
isomorphism.  For $u,v\in H^2(B\mathcal C;k)$, restriction
sends $uv$ to the product of their restrictions on
$B\mathcal D$.  Its vanishing and injectivity in degree four
imply $uv=0$.  No assertion that restriction in degree two is
surjective, or that either inclusion is a homotopy equivalence,
is used.
\end{proof}

\begin{proposition}[Proved: generated long products from a partial map]
\label{A7:generated-partial-long-graph}
Fix positive adjacent hom-sets $A,B,C,D$ and short hom-sets
$U,V,W$, with short maps
\[
 P:A\times B\longrightarrow U,\qquad
 Q:B\times C\longrightarrow V,\qquad
 R:C\times D\longrightarrow W.
\]
Suppose $P$ is completely specified and some entries of $Q$
are specified. Put $U_P=\operatorname{im}P$, and let $V_*$
be the set of values used by the currently specified entries
of $Q$.  Form the graph on the tagged vertex set
\[
 (U_P\times C)\ \sqcup\ (A\times V_*),
\]
inserting, for each specified entry $Q(b,c)$ and every $a\in A$,
the edge
\[
 (P(a,b),c)\;--\;(a,Q(b,c)).
\]
If this graph is connected, every associative category
completing these short products has zero degree-two products
over every field.

Dually, if $R$ is completely specified, put
$W_R=\operatorname{im}R$ and use the graph on
\[
 (B\times W_R)\ \sqcup\ (V_*\times D),
\]
with edges $(b,R(c,d))\;--\;(Q(b,c),d)$ for specified
$Q(b,c)$ and every $d\in D$.  Connectivity gives the same
conclusion. There is no requirement that the specified
entries of $Q$ have already used every label in $V$.
\end{proposition}
\begin{proof}
We prove the source statement.  As another entry of $Q$ is
assigned, an already used value adds only edges.  If the value
$v$ is new, it adds the vertices $(a,v)$ for all $a\in A$.
Each new vertex is immediately joined to the existing vertex
$(P(a,b),c)$ belonging to that new entry.  Hence adding this
entry cannot increase the number of connected components.
The base vertex set $U_P\times C$ is nonempty and is fixed
throughout this process, because $P$ is complete.  A connected
initial graph therefore remains connected after all entries
of $Q$ are assigned.

Fix any completed associative category $\mathcal C$, and let
$\mathcal D$ be its subcategory generated by all adjacent
arrows. Its $02$ and $13$ hom-sets are precisely
$U_P=\operatorname{im}P$ and $V_Q=\operatorname{im}Q$.
Each vertex $(u,c)\in U_P\times C$ gives a composite
$c\circ u:0\to3$ in $\mathcal D$, and each vertex
$(a,v)\in A\times V_Q$ gives $v\circ a$.
Every arrow of $\mathcal D(0,3)$ is obtained in this way.
Associativity identifies the composites at the two ends
of every edge. Since the completed graph is connected,
all these composites coincide, and
$|\mathcal D(0,3)|=1$.
Proposition~\ref{A7:extreme-arrow-collapse}, including its
final consequence, now makes every degree-two product on
$B\mathcal D$ zero over every field.  Lemma
\ref{A7:adjacent-generated-top-homology} transfers this
vanishing to $B\mathcal C$.
Reversing the object order proves the sink statement.
\end{proof}

\begin{remark*}[Scope of the new prefix test]
These graphs deliberately have smaller vertex sets than the
graphs in Corollary~\ref{A7:partial-long-graph}.  Their validity
comes from the actual adjacent-generated subcategory and from
the proved attachment rule for each newly used short value.
One may not simply delete unused vertices from the earlier
universal-category argument.  Completeness of $P$ in the source
test, respectively of $R$ in the sink test, is essential here:
it fixes the nonempty base vertex set to which each later
vertex attaches.

In an initial-prefix search, this certificate excludes every
raw short-map completion of the recorded prefix and every
associative choice of longer products.  Its raw completion
count remains $|U|^{r_P}|V|^{r_Q}|W|^{r_R}$, where $r_P,r_Q,r_R$
are the numbers of unassigned cells.  This counts maps into
the original fixed hom-sets, including all possible future
uses of unused labels; it is not a count of categories or
of normalized representatives.  A finite exclusion additionally
requires an independent prefix check and the complete stream
and profile joins.
\end{remark*}
